%% ======================================================================
%% appA/appendix_a_main.tex
%%
%% Appendix A: Implementation Challenge Narratives
%%
%% Expanded narratives for the nine technical challenges summarized in
%% Table 5.3. Each section follows a Challenge / Response / Lesson
%% structure. The narratives cover: the pivot from the frozen RMRK 2.0
%% pallets to direct use of pallet-nfts; the pivot from a contract-first
%% to a pallet-first architecture; four Snowbridge and XCM integration
%% diagnostics; two Phase-5 security/scaling findings; and the January
%% 2026 ink! discontinuation and its implications.
%%
%% Sections:
%%   A.1 RMRK 2.0 to pallet-nfts pivot
%%   A.2 Contract-First to Pallet-First pivot
%%   A.3 Snowbridge fork-version configuration
%%   A.4 Lodestar preset version pinning
%%   A.5 XCM sovereign account endianness bug
%%   A.6 XCM fee scaling for proof size
%%   A.7 Finding B: authorization gap in xcm_transfer_ownership
%%   A.8 Auto-renewal MaxChildrenReached boundary (50-child cap)
%%   A.9 ink! discontinuation and its implications
%%
%% External labels referenced:
%%   tab:challenges (Table 5.3 in Chapter 5)
%%
%% Internal cross-references:
%%   A.2 (sec:contract-first-pivot) <-> A.9 (sec:ink-discontinuation)
%%
%% Citations: none.
%% ======================================================================

\chapter{Implementation Challenge Narratives}\label{app:challenges}

The nine challenges summarized in Table~\ref{tab:challenges} are presented here in full detail with Challenge, Response, and Lesson for each.

\section{The Pivot from RMRK 2.0 to \texttt{pallet-nfts}}\label{sec:rmrk-to-pallet-nfts}

\textbf{Challenge.} The original Phase~3 design envisioned using the RMRK 2.0 nested-NFT specification, implemented atop the \texttt{rmrk-substrate} pallet suite. RMRK 2.0 precisely offers the parent-child nesting semantics that the CCRMS unified rights record model requires. Reusing an established standard is intended to reduce implementation risk. However, during initial investigations in Phase~4, we discovered that the \texttt{rmrk-substrate} pallets were frozen on Polkadot SDK version 0.9.36 and relied on the now-deprecated \texttt{pallet-uniques} instead of the modern \texttt{pallet-nfts}. Our efforts to compile \texttt{rmrk-substrate} against the current Polkadot SDK monorepo were unsuccessful because of incompatible storage schemas and removed traits.

\textbf{Response.} We decided to replicate RMRK 2.0's nesting semantics directly on top of \texttt{pallet-nfts} without depending on any RMRK pallets. This involved designing a parent-child storage layout (the dual \texttt{Children} and \texttt{Parent} index), implementing the lifecycle helpers (\texttt{mint\_and\_nest\_child}, the burn and re-mint patterns in \texttt{transfer\_ownership}), and tagging child NFTs with a \texttt{rights\_type} attribute via \texttt{pallet\_nfts::set\_attribute}. The implementation took longer than reusing \texttt{rmrk-substrate} would have, but the result is fully integrated with the modern Polkadot SDK and depends only on actively maintained pallets.

\textbf{Lesson.} Established libraries pose latent compatibility risks when they are not kept up to date with the current ecosystem. Reimplementing essential primitives from scratch was justified by the maintainability gain and by the opportunity to demonstrate that RMRK-style composability does not require RMRK itself.

\section[Contract-First to Pallet-First Pivot]{The Pivot from Contract-First to Pallet-First\texorpdfstring{\\}{ }Architecture}\label{sec:contract-first-pivot}

\textbf{Challenge.} The Phase~3 design designated ink! smart contracts as the primary interface for users, featuring three contracts (each corresponding to a different monetization model) to manage user operations. A chain extension mechanism was incorporated to facilitate access to a lightweight pallet dedicated to state coordination. During Phase~4 implementation, two issues arose concurrently: first, the chain extension mechanism within the current \texttt{pallet-revive} release was effectively non-operational for write activities; second, the operational complexity of debugging cross-pallet contract calls was significantly greater than that of direct FRAME pallet interactions.

\textbf{Response.} We inverted the architecture. The pallet (\texttt{pallet-content-rights}) became the definitive source of truth for all business logic, whereas the ink! contract (\texttt{rights\_manager}) was retained as a streamlined user-facing interface that mirrors the pallet's API but holds no independent storage. We developed a custom EVM precompile (\texttt{ContentRightsPrecompile}) to provide read-only access to pallet state for Solidity-based clients, replacing the original chain extension design.

\textbf{Lesson.} The pallet-first decision proved well-founded. The pivot was driven by engineering constraints, and it was taken against the backdrop of the January 2026 announcement that ink! development would be discontinued, which removed any expectation that the contract layer would mature. The episode illustrates that architectural decisions reducing dependence on any single non-essential component become more valuable as the ecosystem evolves.

\section{The Snowbridge Fork-Version Configuration Issue}\label{sec:snowbridge-fork-version}

\textbf{Challenge.} During the integration of Snowbridge, the beacon relay generated Merkle inclusion proofs at incorrect tree positions (gindex 54 instead of 86 for the Electra fork), leading to the rejection of all submitted updates by the Bridge Hub's \texttt{EthereumBeaconClient} pallet.

\textbf{Response.} We identified the root cause as a misinterpretation of the relay's \texttt{forkVersions} configuration field. The relay code interprets this field as a fork activation epoch rather than a version identifier. Consequently, \texttt{0x04000000} was interpreted as the integer 83,886,080, placing the active fork many epochs in the future. We resolved the issue by updating the configuration to use plain epoch numbers (\texttt{\{ "deneb": 0, "electra": 0, "fulu": 5000000 \}}).

\textbf{Lesson.} Cryptographic verification systems are sensitive to silent configuration discrepancies: the failure mode is ``verification fails'' with no informative error pointing to the bad value.

\section{The Lodestar Preset Version Pinning}\label{sec:lodestar-pinning}

\textbf{Challenge.} The \texttt{EthereumBeaconClient} pallet within Bridge Hub is compiled with \texttt{pubkeys: [PublicKey; 512]}. Any beacon header update originating from a smaller sync committee (e.g., the 32-validator minimal preset used by Lodestar's development mode) is rejected during deserialization. Configuring Lodestar v1.41.0 with \texttt{LODESTAR\_PRESET=mainnet} proved ineffective; the development mode silently disregarded the environment variable.

\textbf{Response.} Lodestar v1.35.0, when built from source, adheres to the \texttt{LODESTAR\_PRESET=mainnet} variable. We modified the setup scripts to clone, build, and pin v1.35.0.

\textbf{Lesson.} Cryptographic light clients impose strict requirements that cannot be circumvented through configuration alone. Toolchain compatibility for a research-stage bridge is fragile and must be reconfirmed with each component upgrade.

\section{The XCM Sovereign Account Endianness Bug}\label{sec:xcm-endianness}

\textbf{Challenge.} Our initial cross-chain XCM tests revealed a perplexing failure pattern: the source parachain reported successful XCM dispatch and balance deduction, yet the target parachain indicated \texttt{InsufficientPayment}.

\textbf{Response.} The source chain constructed the sovereign account using polkadot.js's \texttt{toHex()} in default big-endian mode, whereas the CCRMS pallet expected little-endian byte order (the \texttt{SiblingParachainConvertsVia} convention). This produced two separate accounts: one funded but unused, one used but unfunded. We resolved the issue by changing to \texttt{toHex(true)}.

\textbf{Lesson.} Endianness bugs are particularly insidious because both orderings produce valid 32-byte addresses and the symptom (insufficient balance) misdirects the investigation.

\section{The XCM Fee Scaling for Proof Size}\label{sec:xcm-fee-scaling}

\textbf{Challenge.} Initial XCM tests configured the \texttt{BuyExecution} fee at 1B tokens. All tests failed; only at $\sim$70B tokens did they begin to pass.

\textbf{Response.} The dominant cost in the runtime's \texttt{BlockRatioFee} model is \texttt{proof\_size}, not \texttt{ref\_time}. The standard configuration imposes significantly steeper scaling for proof size, meaning even a minimal cross-chain message requires 75--100B tokens. We updated test scripts to specify 100B tokens.

\textbf{Lesson.} Polkadot fee modeling is driven primarily by proof size, not computational effort.

\section[Finding B: Authorization Gap]{Finding B: The Authorization Gap in\texorpdfstring{\\}{ }\texttt{xcm\_transfer\_ownership}}\label{sec:xcm-auth-gap}

\textbf{Challenge.} During the Phase~5 security audit, we identified an authorization flaw in \texttt{xcm\_transfer\_ownership}: the initial implementation performed only \texttt{ensure\_signed(origin)} before burning/minting NFTs based on a \texttt{from} parameter, allowing any signed account to steal ownership.

\textbf{Response.} We added \texttt{ensure!(authorizer == from, Error::<T>::Unauthorized)} and committed a unit test (\texttt{security\_xcm\_transfer\_ownership\_any\_account\_can\_steal}) that verifies the exploit is blocked.

\textbf{Lesson.} Any cross-chain extrinsic accepting a ``from'' identity parameter requires explicit authorization. Payer-beneficiary separation introduces impersonation risk absent from local-only calls.

\section[Auto-Renewal Child Cap]{The Auto-Renewal \texttt{MaxChildrenReached}\texorpdfstring{\\}{ }Boundary}\label{sec:children-bound}

\textbf{Challenge.} Batches of 100 \texttt{subscribe} calls to the same content item failed midway with \texttt{MaxChildrenReached}. The \texttt{Children} storage entry is limited to 50 child NFTs by \texttt{BoundedVec<\_, ConstU32<50>{}>}.

\textbf{Response.} We documented the constraint as a known limitation and updated benchmark scripts to rotate content items. Subsequent analysis showed the bound to be stricter than a concurrency limit. Children are minted per buyer and released only when a view pack is exhausted, so expired subscriptions permanently occupy their slots. The bound therefore limits a content item to fifty subscribers and owners over its lifetime. Its implications for the economic results are discussed in Section~\ref{sec:limitations}.

\textbf{Lesson.} Bounded FRAME collections trade benchmarkability for capacity limits, and a bound that suits a benchmark may conflict with the target workload. Bounds should be checked against the audience sizes the system is meant to serve, and redundant indexes should be removed rather than bounded.

\section[ink! Discontinuation]{The ink! Discontinuation and Its Implications}\label{sec:ink-discontinuation}

\textbf{Challenge.} On 27~January 2026, the ink! Alliance announced the cessation of ink! development following three failed OpenGov treasury funding proposals. Both the Web3 Foundation and Parity Technologies declined continued funding. The announcement outlined three migration options: Solidity via \texttt{revive}, low-level Rust against \texttt{pallet-revive}, or the community-maintained \texttt{wrevive} SDK.

\textbf{Response.} Because our pallet-first architecture (Appendix~\ref{sec:contract-first-pivot}) centers business logic in a FRAME pallet rather than in ink! contracts, the discontinuation has minimal practical impact. The ink! code in \texttt{contracts/rights\_manager} serves as a non-load-bearing wrapper. We deferred migration.

\textbf{Lesson.} Architectural decisions that minimize dependence on any single non-essential component provide advantages during ecosystem-wide changes.
